\section{Worked benchmark: the stable cubic \(y^3+y=x\)}

This section gives an end-to-end mathematical example suitable for the first
full object-language implementation of the framework.  It uses only exact
rational arithmetic, bounded bisection, and quantifier-free induction.

Fix a rational input \(x\) satisfying \(0\le x\le1\), and set
\[
       F_x(y)=y^3+y-x.
\]
We construct a regular name for the unique solution in \([0,1]\).

\subsection{Exact bisection}

Set \(\ell_0=0\) and \(u_0=1\).  Given \(\ell_n,u_n\), let
\[
      m_n=\frac{\ell_n+u_n}{2}.
\]
Since \(F_x(m_n)\) is rational, its sign is decidable exactly.  Define
\[
(\ell_{n+1},u_{n+1})=
\begin{cases}
 (m_n,u_n),&F_x(m_n)\le0,\\
 (\ell_n,m_n),&F_x(m_n)>0.
\end{cases}
\]

\begin{lemma}[Bisection invariant]
For every \(n\),
\[
  0\le\ell_n\le u_n\le1,\qquad
  F_x(\ell_n)\le0\le F_x(u_n),\qquad
  u_n-\ell_n=2^{-n}.
\]
\end{lemma}

\begin{proof}
At \(n=0\), \(F_x(0)=-x\le0\), while
\(F_x(1)=2-x\ge1\).  Suppose the invariant holds.  The midpoint lies between
the endpoints.  If \(F_x(m_n)\le0\), the left endpoint can be replaced by
\(m_n\); otherwise the right endpoint can be replaced by \(m_n\), and exact
rational order gives \(0<F_x(m_n)\).  In either case the sign bracket is
preserved and the width is halved.  The entire invariant is a conjunction of
rational cross-product inequalities and one dyadic-width equality, so its
induction matrix is quantifier-free.
\end{proof}

Let
\[
        c_n=\frac{\ell_n+u_n}{2}.
\]

\begin{lemma}[Regularity]
If \(m\ge n\), then
\[
          \abs{c_m-c_n}\le 2^{-n}.
\]
Consequently \(a(k)=c_{k+2}\) is a fast regular name.
\end{lemma}

\begin{proof}
The bisection intervals are nested.  A finite induction on \(m-n\) proves
\[
       [\ell_m,u_m]\subseteq[\ell_n,u_n].
\]
Both centers lie in the latter interval, whose width is \(2^{-n}\).  For the
shifted name,
\[
 \abs{a(j)-a(k)}
 \le 2^{-(k+2)}<2^{-k}
 \qquad(j\ge k).
\]
No completeness axiom is used in the internal statement; regularity is a
family of exact rational inequalities.
\end{proof}

\subsection{Residual bound}

For \(0\le r\le s\le1\),
\[
\begin{aligned}
  F_x(s)-F_x(r)
  &=(s-r)(s^2+sr+r^2+1),\\
  s-r
  &\le F_x(s)-F_x(r)\le4(s-r).
\end{aligned}
\]
The upper bound uses \(s^2,sr,r^2\le1\).

\begin{theorem}[Certified cubic name]
For every \(n\),
\[
       \abs{F_x(c_n)}\le 4\cdot2^{-n}.
\]
In particular, \(a(k)=c_{k+3}\) satisfies
\[
       \abs{a(k)^3+a(k)-x}\le2^{-k-1}.
\]
\end{theorem}

\begin{proof}
Monotonicity places \(F_x(c_n)\) between
\(F_x(\ell_n)\le0\) and \(F_x(u_n)\ge0\).  Hence
\[
 \abs{F_x(c_n)}
 \le F_x(u_n)-F_x(\ell_n)
 \le4(u_n-\ell_n)=4\cdot2^{-n}.
\]
All quantities are rational, and the factorization of the cubic difference
is a polynomial identity.
\end{proof}

\subsection{Quantitative uniqueness}

\begin{lemma}[Residual isolation]
For rational \(p,q\in[0,1]\),
\[
       \abs{p-q}\le \abs{F_x(p)-F_x(q)}
       \le \abs{F_x(p)}+\abs{F_x(q)}.
\]
Thus if both residuals are at most \(\eps\), then
\(\abs{p-q}\le2\eps\).
\end{lemma}

\begin{proof}
The factorization gives
\[
 \abs{F_x(p)-F_x(q)}
 =\abs{p-q}\,(p^2+pq+q^2+1).
\]
The second factor is at least \(1\), proving the first inequality.  The
second is the triangle inequality.
\end{proof}

It follows that any two regular names in \([0,1]\) whose residuals tend to
zero with explicit moduli are extensionally equal.  The proof is quantitative:
to establish distance at most \(2^{-k}\), request residuals at most
\(2^{-k-2}\) and approximations accurate enough that their rounding errors
consume the remaining half of the budget.

\subsection{The exact certificate}

A finite certificate for request \(k\) contains:
\[
  x,\quad n=k+3,\quad
  (\ell_0,u_0),\ldots,(\ell_n,u_n),
\]
together with the branch bit at each stage.  A verifier checks:
\begin{enumerate}[label=(\roman*)]
\item \(x\) and every endpoint are valid rational codes;
\item \(0\le x\le1\);
\item the initial interval is \([0,1]\);
\item each successor is the prescribed half selected by the recorded exact
      sign test;
\item every sign bracket and width equality holds;
\item the final center and residual bound are correct.
\end{enumerate}
The trace length is \(k+3\), a primitive-recursive function of the free
request.  The invariant is precisely the bisection invariant above.  This is
therefore a complete instance of the finite-trace reflection theorem.

\subsection{Extension to represented \(x\)}

If \(x\) is itself given by a regular name, exact midpoint signs are
unavailable near zero.  Use the overlapping-branch rule:
at stage \(n\), approximate \(F_x(m_n)\) within \(\eta_n\).  A confidently
negative value selects the right half, a confidently positive value selects
the left half, and an ambiguous value returns a short central interval whose
radius accounts for \(\eta_n\) using the lower slope bound \(1\).
The output is initially relational; the residual isolation lemma proves that
all accepted intervals are mutually close.  Reverse demand propagation picks
the input precision needed to make the final interval width and residual
smaller than \(2^{-k}\).  This demonstrates why the relational framework is
not a concession: it is the stable uniform form of bisection on represented
inputs.
